% ============================================================
%  capitulos/cap1.tex  —  AUTECH · Trabajo Terminal ESCOM-IPN
%  Capítulo 1: Marco introductorio
% ============================================================

\chapter{Marco introductorio}
\label{ch:cap1}

% ─────────────────────────────────────────────────────────────
\section{Introducción}
\label{sec:introduccion}

La Teoría de la Computación es una asignatura de gran importancia dentro del plan de estudios
de Ingeniería en Sistemas Computacionales de la ESCOM, ya que establece las bases matemáticas
y lógicas que sustentan áreas como el diseño de compiladores, la arquitectura de computadoras
y el desarrollo de lenguajes de programación~\cite{ipn2020, aho2006}. Sin embargo, la
naturaleza abstracta de sus temas como autómatas finitos, lenguajes formales y Máquinas de
Turing convierte su enseñanza en un reto constante para docentes y alumnos por
igual~\cite{sinchi2018, felder2016}.

Este problema no es exclusivo de la ESCOM. En general, las materias de alto contenido formal
dentro de las carreras STEM presentan altos índices de dificultad de comprensión cuando el
estudiante no cuenta con herramientas que le permitan experimentar con los conceptos, más allá
del pizarrón y los apuntes~\cite{munoz2012}. La práctica continua y el contacto directo con
los modelos son factores clave para consolidar este tipo de
conocimientos~\cite{hopcroft2007}.

Frente a esto, el aprendizaje móvil (m-learning) ha demostrado ser una alternativa eficaz para
fortalecer el aprendizaje de conceptos abstractos al estudiante de forma visual e interactiva,
aprovechando la disponibilidad que ofrece un dispositivo móvil~\cite{crompton2018}. Bajo esa
premisa nace AUTECH: una plataforma educativa para Android orientada a fortalecer el aprendizaje
de la Teoría de la Computación en la ESCOM, integrando un motor de simulación de autómatas, un
sistema de aprendizaje por niveles y componentes de ludificación que buscan mantener la
motivación del usuario a lo largo del proceso, además de un sistema de gestión para que tanto
alumnos como profesores puedan trabajar en conjunto~\cite{deterding2011}.

Para guiar al lector a través del proceso de ingeniería seguido en este proyecto, el documento
se organiza de la siguiente manera:

\begin{itemize}
    \item \textbf{Capítulo 1 -- Marco Introductorio:} Establece el contexto académico, el
        planteamiento del problema, los objetivos y las fronteras del
        proyecto~\cite{ipn2020, escom2016}.

    \item \textbf{Capítulo 2 -- Marco Teórico y Estado del Arte:} Presenta los fundamentos de
        la Teoría de la Computación, los principios de diseño UI/UX aplicados al m-learning y
        una comparativa con las herramientas existentes~\cite{sinchi2018, crompton2018}.

    \item \textbf{Capítulo 3 -- Análisis:} Documenta el levantamiento de requerimientos bajo
        el estándar IEEE 830, el modelado funcional en UML, el análisis de riesgos, la
        factibilidad del proyecto y los términos legales y de privacidad
        aplicables~\cite{ieee1998}.

    \item \textbf{Capítulo 4 -- Diseño:} Detalla las decisiones de arquitectura, el diseño
        de datos, los mockups de interfaz, los diagramas del sistema y la infraestructura
        tecnológica, integrando Design Thinking como enfoque centrado en el usuario.

    \item \textbf{Capítulo 5 -- Ciclo~0, Prototipo Base (MVP):} Describe la implementación
        del primer ciclo de desarrollo evolutivo, abarcando el motor de simulación de autómatas
        finitos y los módulos de aprendizaje básicos.
\end{itemize}

Los capítulos correspondientes a los ciclos 1 y 2 del prototipado evolutivo serán incorporados
durante la fase TT2, conforme al plan de desarrollo establecido en la
sección~\ref{ssec:ciclos}.

% ─────────────────────────────────────────────────────────────
\section{Antecedentes}
\label{sec:antecedentes}

Históricamente, la enseñanza de la Teoría de la Computación en la ESCOM se ha apoyado en el
desarrollo analítico de modelos matemáticos mediante representaciones gráficas tradicionales
como diagramas en pizarrón, ejercicios en papel y sesiones de laboratorio
presenciales~\cite{ipn2020}. Con el tiempo, la academia incorporó simuladores de escritorio
para hacer estos conceptos más tangibles, siendo JFLAP la herramienta que se escogió como
estándar en las prácticas de la institución~\cite{rodger2006}.

Si bien JFLAP representó un avance significativo en su momento, actualmente las condiciones
del entorno educativo han cambiado. Las herramientas de escritorio, aunque siguen siendo
funcionales, están atadas a equipos fijos y no han evolucionado al ritmo de las necesidades
actuales del estudiante, quien cada vez más organiza su estudio fuera del
laboratorio~\cite{nielsen1994}. Frente a esto, la oferta de aplicaciones móviles relacionadas
con ToC es escasa y, cuando existe, se limita a simuladores técnicos sin estructura didáctica
o una estructura didáctica sin una pizca de ludificación, lo que no contribuye a una buena
comprensión o al menos limita la comprensión de modelos a un uso mecánico de los mismos,
además de que tampoco contribuye a un alto índice de retención en los sistemas
actuales~\cite{felder2016, crompton2018}.

Este contexto evidencia una brecha entre las herramientas disponibles y las formas en que el
estudiante actual aprende y se organiza, brecha que AUTECH busca atender.

% ─────────────────────────────────────────────────────────────
\section{Planteamiento del Problema}
\label{sec:planteamiento}

En la ESCOM, el aprendizaje de la Teoría de la Computación enfrenta una barrera que va más
allá de la dificultad natural de la materia: la alta abstracción de sus modelos y temas genera
una brecha con la familiarización que los métodos de enseñanza actuales no logran cerrar del
todo~\cite{hopcroft2007}. Aunque los índices de acreditación de la asignatura se perciben como
aceptables, se observa de manera empírica un fenómeno de arrastre de deficiencias conceptuales
que se hace evidente en materias siguientes del plan de estudios ISC
2020~\cite{ipn2020, felder2016}.

El problema concreto es que muchos estudiantes desarrollan una comprensión mecánica suficiente
para aprobar evaluaciones inmediatas, pero sin la asimilación profunda necesaria para aplicar
esos conocimientos en contextos reales, como el diseño de analizadores léxicos o la
implementación de lógica de control en hardware~\cite{aho2006}. Esto se llega a convertir en
un cuello de botella cuando llegan a materias como diseño de compiladores o arquitectura de
computadoras, que dan por sentado un dominio sólido de estos fundamentos~\cite{munoz2012}.

Parte del problema radica en que las herramientas de apoyo actuales no favorecen la práctica
continua. Al estar limitadas a entornos de escritorio, el contacto del alumno con los modelos
queda limitado a sesiones específicas de laboratorio, lo cual no es suficiente para unificar
conocimiento de esta naturaleza~\cite{escom2016, nielsen1994}.

% ─────────────────────────────────────────────────────────────
\section{Solución Propuesta}
\label{sec:solucion}

Como respuesta a la problemática identificada, se propone el desarrollo de AUTECH, una
plataforma educativa móvil para Android diseñada para fortalecer el aprendizaje de la Teoría
de la Computación en la ESCOM mediante simulación interactiva, aprendizaje guiado y
ludificación~\cite{deterding2011}. El contenido se organiza en niveles de progresión alineados
con el programa sintético de la asignatura, siguiendo la estructura formal de la Jerarquía de
Chomsky~\cite{church1936} y los fundamentos computacionales establecidos por
Turing~\cite{turing1936}, desde autómatas finitos hasta Máquinas de Turing, exceptuando
los lenguajes tipo~1 ya que no están contemplados en el programa sintético de
ISC 2020~\cite{ipn2020, hopcroft2007}.

\begin{infobox}
\textbf{Nota:} Aunque la aplicación AUTECH está desarrollada en lenguaje Dart, con ayuda del
framework multiplataforma Flutter, el alcance de este Trabajo Terminal restringe el despliegue
y distribución exclusivamente al sistema operativo Android. Esta delimitación técnica resulta
de tres factores críticos: primero, \textbf{infraestructura de compilación}: la generación de
ejecutables para iOS exige hardware propietario (macOS con Xcode), del cual carecemos y no
podemos costear; segundo, \textbf{sostenibilidad financiera}: el ecosistema de Apple requiere
un licenciamiento anual de \$99~dólares, lo cual es inviable para un proyecto académico sin
fines de lucro; y tercero, \textbf{agilidad de distribución}: Android permite el despliegue
directo de paquetes instalables (\textit{sideloading} mediante .apk), lo que resulta vital
para realizar pruebas de usuario rápidas bajo la metodología de Prototipado Evolutivo. Además,
estadísticamente, priorizar Android asegura la compatibilidad con más del 75\% del mercado
móvil estudiantil en México según el reporte de StatCounter global stats realizado en abril
del 2026.
\end{infobox}

La plataforma se articula en tres componentes principales:

\textbf{Motor de Simulación y Validación Interactiva.} Desarrollado con Flutter y el motor
Flame, permite al estudiante construir y manipular autómatas directamente en la pantalla
mediante gestos táctiles. El motor valida los modelos probando cadenas de entrada y soporta
conversiones automáticas entre modelos equivalentes, como la transformación de un AFND a un
AFD, facilitando la experimentación y comprensión de las relaciones entre
modelos~\cite{sinchi2018, rodger2006}.

\textbf{Sistema de Gestión de Aprendizaje y Analíticas.} AUTECH opera bajo un esquema
multi-rol (Administrador, Profesor y Alumno) respaldado por Supabase como infraestructura de
base de datos relacional en la nube, ya que, aunque anteriormente se tenía pensado el uso de
Firebase, se optó por el cambio debido a los costos que implicaba pagar Firebase para lograr
implementar el envío de correos de confirmación. El docente cuenta con un panel de analíticas
que visualiza el desempeño grupal e individual, permitiendo identificar temas problemáticos y
asignar actividades de refuerzo a grupos específicos de alumnos, además de estadísticas de
temas o ejercicios con más problemas de asimilación por parte del
alumno~\cite{felder2016, munoz2012}.

\newpage
\textbf{Experiencia de Usuario Ludificada.} El recorrido de aprendizaje combina teoría y
práctica en cada lección, integrando ejercicios de opción múltiple, construcción de autómatas
y depuración de modelos incorrectos. Un sistema de insignias y rachas incentiva la constancia
del usuario, mientras que la persistencia local de datos garantiza el funcionamiento de la
aplicación en entornos sin conexión~\cite{deterding2011, materialdesign2023}.

La arquitectura del sistema sigue una arquitectura por capas basada en Layered Architecture
con estructura modular por funcionalidad, anteriormente concebida como MVC en el protocolo
inicial del proyecto. La transición respondió a una decisión técnica surgida durante el primer
ciclo de desarrollo: al implementar el sistema en Flutter, la naturaleza del framework favorece
que las pantallas consuman los servicios directamente mediante \texttt{StreamBuilder} y
\texttt{FutureBuilder}, sin requerir una capa de controlador explícita. Forzar una separación
MVC estricta habría introducido complejidad innecesaria sin beneficio real para el tamaño y
alcance del proyecto. La arquitectura adoptada mantiene la separación de responsabilidades
mediante dos capas bien definidas: una capa de datos compuesta por servicios y modelos, y una
capa de UI organizada por módulos funcionales, lo cual es coherente con los principios de
mantenibilidad y escalabilidad que originalmente motivaron la elección de
MVC~\cite{pressman2010, reenskaug1979}.

El aporte central de ingeniería de AUTECH no reside en la gestión convencional de datos, sino
en la integración de cuatro elementos técnicos no triviales en una sola plataforma móvil: un
motor de simulación algorítmica capaz de ejecutar y visualizar en tiempo real modelos
computacionales formales, incluyendo la construcción de subconjuntos para conversión AFND a
AFD, la minimización de estados por el algoritmo de Hopcroft, la construcción de Thompson para
expresiones regulares y la simulación determinista de Máquinas de Turing; un patrón de
resiliencia API-First con Fallback Local que garantiza disponibilidad total independientemente
de la conectividad del dispositivo (exceptuando el material audiovisual que será externo a la
app); y un sistema de analíticas educativas que transforma el comportamiento del estudiante
dentro de la plataforma en métricas accionables para el docente. Además, el módulo de
aprendizaje guiado utiliza componentes de ludificación para obtener mayores índices de
retención por parte de los alumnos dentro de la aplicación. La combinación de estos cuatro
elementos, orientada específicamente al contexto académico de la ESCOM y alineada con el
programa sintético ISC 2020, constituye la contribución técnica distintiva del
proyecto~\cite{hopcroft2007, pressman2020, flame2024}.

\newpage
% ─────────────────────────────────────────────────────────────
\section{Metodología de Desarrollo: Paradigma de Prototipado evolutivo}
\label{sec:metodologia}

Para la construcción de AUTECH se adoptó el Paradigma de Construcción de Prototipos en su
variante evolutiva, fundamentado en la metodología de Roger~S.\ Pressman~\cite{pressman2010}.
A diferencia del prototipado desechable, en este enfoque cada prototipo no se descarta al
finalizar el ciclo, sino que se convierte en la base del siguiente, permitiendo que el sistema
crezca y se refine de forma progresiva. Esta característica lo hace especialmente adecuado para
el desarrollo de herramientas didácticas móviles, donde los requisitos pedagógicos y técnicos
se validan y ajustan con cada entrega.

Como se ilustra en la Figura~\ref{fig:paradigma_prototipos}, el ciclo de vida del proyecto se
compone de cinco etapas que se repiten en cada iteración, además de visualizarlos en la
Ilustración~\ref{fig:paradigma_prototipos}:

\begin{itemize}
    \item \textbf{Comunicación:} Levantamiento de requerimientos a partir de las necesidades
        de aprendizaje identificadas en la ESCOM para los temas de teoría de la computación.

    \item \textbf{Plan rápido:} Definición del alcance de cada entrega y gestión de los
        riesgos técnicos asociados~\cite{pmi2021}.

    \item \textbf{Modelado y diseño rápido:} Traducción de los requerimientos en diagramas
        UML y maquetas de alta fidelidad en Figma, estableciendo la experiencia de usuario
        antes de escribir código.

    \item \textbf{Construcción del prototipo:} Implementación del código fuente con Dart,
        Flutter y Flame Engine, produciendo un simulador funcional de autómatas como entregable
        de cada ciclo.

    \item \textbf{Despliegue, entrega y retroalimentación:} El prototipo se entrega a
        estudiantes y docentes para pruebas de usabilidad, esto a partir del ciclo~1. Los
        comentarios obtenidos reinician el ciclo, orientando las mejoras de la siguiente
        iteración~\cite{pressman2010}.
\end{itemize}

\begin{figure}[H]
    \centering
    \includegraphics[width=0.70\textwidth]{fig_paradigma_prototipos}
    \caption{Paradigma de Construcción de Prototipos.}
    \label{fig:paradigma_prototipos}
\end{figure}

% ──────────────────────────────────────────────────────────
\subsection{Justificación del Enfoque Evolutivo}
\label{ssec:justificacion_evolutivo}

La elección del prototipado evolutivo sobre modelos secuenciales como el modelo en cascada
responde a la naturaleza del proyecto: un simulador interactivo con alta complejidad
algorítmica y visual, cuyos requisitos didácticos solo pueden validarse con usuarios
reales~\cite{pressman2020}. Las ventajas concretas para AUTECH son:

\textbf{Reducción de riesgos técnicos.} Validar el motor gráfico y la lógica de validación
de cadenas en modelos simples como los autómatas finitos antes de abordar arquitecturas más
complejas como las Máquinas de Turing reduce significativamente la probabilidad de errores
estructurales en etapas avanzadas~\cite{aho2006, pressman2020}.

\textbf{Adaptabilidad didáctica.} Evaluar la retención y usabilidad del estudiante en las
primeras etapas permite calibrar la dificultad de los niveles y los componentes de
ludificación con datos reales, en lugar de suposiciones
iniciales~\cite{deterding2011, nielsen1994}.

\textbf{Entrega temprana de valor.} Cada ciclo produce un módulo funcional que puede probarse
en el entorno académico real, sin necesidad de esperar a que el sistema esté
completo~\cite{escom2016}.

Este principio se evidencia en la propia evolución del proyecto: durante el primer ciclo de
desarrollo, la arquitectura y el stack tecnológico definidos en el protocolo inicial fueron
refinados como resultado directo de la validación técnica. El prototipado evolutivo no solo
permite, sino que anticipa este tipo de ajustes: cada ciclo aporta información real que orienta
las decisiones del siguiente, haciendo del cambio controlado una fortaleza del paradigma, sin
desviarnos completamente de él~\cite{pressman2010}.

% ──────────────────────────────────────────────────────────
\subsection{Definición de los Ciclos de Iteración}
\label{ssec:ciclos}

El proyecto se estructura en tres ciclos de desarrollo, cada uno escalando progresivamente en
complejidad teórica y funcionalidad operativa:

\textbf{Ciclo~0: Núcleo y Lenguajes Regulares.} Construcción de la arquitectura base bajo un
esquema de Layered Architecture con estructura modular por funcionalidad. Incluye el motor de
simulación para Autómatas Finitos Deterministas y No Deterministas, las validaciones lógicas
de cadenas, además de la simulación de Autómatas de pila y la gestión inicial de usuarios
mediante Supabase~\cite{pressman2020, supabase2024}. En este ciclo solo se incluyen pruebas
internas que no consideran la validación con estudiantes y/o profesores, debido a que aún no
estaría implementado por completo el módulo de gestión y estadísticas.

\textbf{Ciclo~1: Gestión Docente.} Se integra el panel de analíticas en aplicación desktop
para el docente, con visualización semaforizada del desempeño grupal e
individual~\cite{materialdesign2023}. Refinación del prototipo anterior basado en validaciones
con estudiantes y profesores. Y se integra en panel del simulador, la opción de conversión de
un autómata pila a su gramática independiente del contexto equivalente.

\textbf{Ludificación Avanzada:} Consolidación del ecosistema completo. Se despliega el
sistema de insignias y rachas para fortalecer la autonomía y retención del
estudiante~\cite{deterding2011}. Refinación del prototipo anterior basado en validaciones con
estudiantes y profesores, además de los manuales correspondientes.

\textbf{Ciclo~2:} Se integra al simulador la opción de simulación de máquinas de Turing
unicinta, así como las opciones de operaciones sobre lenguajes regulares, conversión de
AFN, AFN-$\varepsilon$ a AFD, y minimización de estados, e integra las notificaciones de
estilo push.

% ──────────────────────────────────────────────────────────
\subsection{Cronograma de actividades}
\label{ssec:cronograma}

\begin{table}[H]
    \centering
    \caption{Cronograma de actividades generales para cada ciclo de desarrollo.}
    \label{tab:cronograma}
    \includegraphics[width=\textwidth]{tabla_cronograma}
\end{table}

% ─────────────────────────────────────────────────────────────
\section{Objetivos}
\label{sec:objetivos}

\subsection{Objetivo General}
\label{ssec:obj_general}

Diseñar e implementar una plataforma educativa móvil ludificada para Android que permita a los
estudiantes de Teoría de la Computación de la ESCOM simular, construir y validar modelos
computacionales de la Jerarquía de Chomsky \cite{church1936}, cuya base teórica formal
fue establecida por Turing \cite{turing1936}, además de contar con un sistema de aprendizaje
guiado por niveles, analíticas de desempeño para el docente y soporte multi-rol.

\subsection{Objetivos Específicos}
\label{ssec:obj_especificos}

\begin{itemize}
    \item Implementar un motor de simulación gráfica interactiva que permita la construcción,
        manipulación táctil y validación de modelos computacionales, incluyendo conversiones
        automáticas entre modelos equivalentes.

    \item Diseñar un sistema de aprendizaje guiado con progresión por niveles, integrando
        teoría, ejercicios prácticos y retroalimentación inmediata para el estudiante.

    \item Desarrollar un panel de a nalíticas en desktop orientado al perfil docente, capaz de
        procesar y visualizar métricas de desempeño grupal e individual mediante indicadores
        semaforizados.

    \item Integrar un mecanismo de persistencia y sincronización de datos en la nube mediante
        Supabase, que garantice el funcionamiento de la aplicación en entornos sin conexión y
        salvaguarde el progreso académico~\cite{supabase2024}.

    \item Construir un módulo de administración global para la gestión centralizada de usuarios
        y control de acceso al sistema.
\end{itemize}

% ─────────────────────────────────────────────────────────────
\section{Justificación}
\label{sec:justificacion}

El desarrollo de AUTECH responde a la necesidad de modernizar las herramientas didácticas
disponibles para la asignatura de Teoría de la Computación en la ESCOM, abordando la
problemática desde tres dimensiones: la eficiencia académica, la accesibilidad tecnológica y
la alineación institucional.

\subsection{Justificación Académica}
\label{ssec:just_academica}

La Teoría de la Computación es base obligada para materias de alto rigor técnico en semestres
superiores, como el diseño de compiladores y la arquitectura de
computadoras~\cite{ipn2020, aho2006}. La asimilación superficial de sus conceptos llega a
generar un fenómeno de arrastre de deficiencias que compromete la trayectoria del estudiante:
al existir una dependencia seriada de conocimientos, la falta de comprensión en esta materia
frena el avance y puede derivar en índices más altos de recurse en ciertas materias, lo que a
su vez en ocasiones llega a saturar la oferta educativa~\cite{tinto1987}.

Un factor que hace que esta problemática sea más difícil de afrontar es que los métodos de
enseñanza tradicionales, centrados en la pizarra y los apuntes, privilegian un único canal de
transmisión del conocimiento. Sin embargo, los estudiantes no aprenden de forma homogénea:
el modelo VARK identifica perfiles de aprendizaje visual, auditivo, lector/escritor y
kinestésico, cuyas necesidades cognitivas difieren significativamente entre
sí~\cite{fleming1992}. En un entorno donde la herramienta principal es JFLAP ---una interfaz
estática operada desde escritorio---, los estudiantes kinestésicos y visuales, que requieren
manipulación directa y representaciones dinámicas para consolidar conceptos abstractos, se
encuentran en desventaja sistemática~\cite{felder2016}.

AUTECH atiende esta brecha mediante un diseño multimodal: la construcción táctil de autómatas
responde al perfil kinestésico, las animaciones y diagramas dinámicos al perfil visual, y el
contenido teórico estructurado por lección al perfil lector. Como trabajo futuro, está
considerada la implementación de elementos de sonido para incluir los perfiles de aprendizaje
auditivos junto con el material audiovisual, sin olvidar los elementos hápticos que pueden ser
considerados de ser necesario después del ciclo~1. Esta convergencia de canales no busca validar
empíricamente los estilos de aprendizaje, sino reducir la dependencia de un único modo de
transmisión del conocimiento y ampliar el alcance didáctico de la
herramienta~\cite{felder2016, sweller1988}.

\subsection{Justificación Tecnológica y de Movilidad}
\label{ssec:just_tecnologica}

Aunque la comunidad estudiantil de la ESCOM tiene acceso a equipos de cómputo, la dependencia
de software de escritorio como JFLAP restringe el estudio a entornos fijos y sesiones de
laboratorio programadas~\cite{rodger2006}. En un contexto donde un sector considerable de los
estudiantes invierte hasta cinco horas diarias en traslados metropolitanos, la portabilidad se
convierte en un factor relacionado con el rendimiento académico~\cite{kukulska2007}. AUTECH
permite convertir esos tiempos muertos en sesiones de práctica activa, aprovechando la
disponibilidad del dispositivo móvil y mitigando además los problemas de conectividad mediante
su modo sin conexión~\cite{crompton2018, materialdesign2023}.

\subsection{Justificación Institucional}
\label{ssec:just_institucional}

El Instituto Politécnico Nacional promueve activamente la incorporación de tecnologías que
modernicen y dinamicen sus procesos educativos a través del modelo educativo del IPN, así como
el documento rector de operación y evaluación para los trabajos terminales. Sin embargo, con
frecuencia existe una brecha entre los objetivos del programa de estudios y las herramientas
realmente disponibles en el aula~\cite{ipn2020, escom2016}. AUTECH se alinea con esa visión
institucional al ofrecer una solución tecnológica original, desarrollada desde la propia
comunidad de la ESCOM, que responde a las necesidades concretas de sus estudiantes y docentes.

% ─────────────────────────────────────────────────────────────
\section{Resultados Esperados}
\label{sec:resultados}

Al concluir el desarrollo e implementación del proyecto, se proyecta la obtención de los
siguientes resultados tangibles y medibles, orientados a satisfacer las necesidades de
actualización tecnológica en la Academia de Ciencias de la Computación de la
ESCOM~\cite{ipn2020, escom2016}.

\subsection{Ecosistema de Software Funcional}
\label{ssec:ecosistema}

\begin{itemize}
    \item \textbf{Aplicación Móvil (AUTECH):} Un paquete de instalación (.APK) para el
        sistema operativo Android que integre el motor de simulación de autómatas, el módulo
        de aprendizaje guiado y el validador algorítmico de lenguajes formales, además de los
        módulos anteriormente mencionados~\cite{crompton2018}.

    \item \textbf{Panel de Gestión en aplicación de escritorio (Dashboard):} Una plataforma
        administrativa para el cuerpo docente que, mediante Supabase como infraestructura de
        base de datos relacional en la nube, permita visualizar de forma centralizada las
        métricas de desempeño grupal e individual de múltiples grupos y/o carreras de la
        ESCOM~\cite{supabase2024}.
\end{itemize}

\subsection{Métricas de Éxito y Validación Didáctica}
\label{ssec:metricas}

El impacto del proyecto no se medirá basándose en índices de aprobación externos, sino a
través de la evaluación directa del software mediante instrumentos estandarizados:

\begin{itemize}
    \item \textbf{Validación de Usabilidad:} Obtención de un puntaje superior al 80\% de
        aceptación en la Escala de Usabilidad del Sistema (SUS) aplicada a una muestra
        representativa de estudiantes activos y consideración de un modelo TAM~\cite{brooke1996}
        para los ciclos~1 y~2.

    \item \textbf{Autopercepción de Competencias:} Confirmación, mediante instrumentos de
        evaluación post-uso, de una reducción en la carga cognitiva y un aumento en la
        confianza técnica de los usuarios al modelar los componentes específicos de la
        Jerarquía de Chomsky~\cite{hopcroft2007, sweller1988}.

    \item \textbf{Analíticas de Intervención:} Generación de un registro histórico (o al
        menos con un historial de 3~semestres) en la nube que mapee los errores lógicos más
        frecuentes por grupo, dotando al docente de datos empíricos para realizar
        intervenciones didácticas de precisión~\cite{felder2016}.
\end{itemize}

\subsection{Sistema de Progresión y Ludificación}
\label{ssec:ludificacion}

El componente de ludificación producirá un historial de competencias digitales para cada
usuario, estructurado en los siguientes elementos:

\begin{itemize}
    \item \textbf{Jerarquía de Insignias:} Un sistema de reconocimiento progresivo con
        insignias alusivas a las secciones completadas que clasifica el dominio del estudiante
        a partir de la resolución de problemas y la depuración de cadenas, además de los
        típicos cuestionarios~\cite{deterding2011}.

    \item \textbf{Desbloqueo Secuencial de Mundos:} Evidencia visual del avance lógico del
        alumno a través de los modelos de la Jerarquía de Chomsky, desde los Autómatas
        Finitos hasta las Máquinas de Turing, con la excepción mencionada anteriormente de
        los lenguajes tipo~1~\cite{hopcroft2007}.

    \item \textbf{Indicador de Racha:} Registro visual del número de días consecutivos en
        que el alumno ha completado ejercicios, fomentando la práctica continua y la
        constancia en el estudio.
\end{itemize}

\subsection{Interoperabilidad y Transferencia Tecnológica}
\label{ssec:interoperabilidad}

\textbf{Compatibilidad con JFLAP y FSAM:} Desarrollo de un módulo de análisis (parser) capaz
de importar y exportar la estructura XML de los archivos .jff, garantizando que los autómatas
construidos en el dispositivo móvil sean bidireccionalmente compatibles con el software de
escritorio utilizado en los laboratorios de la institución~\cite{rodger2006}, además de la
implementación de archivos exportados desde FSAM.

\textbf{Documentación de Ingeniería:} Entrega de un repositorio de código fuente organizado
bajo el patrón de Layered Architecture y manuales correspondientes, asegurando que el sistema
sea mantenible y escalable por futuras generaciones de la ESCOM~\cite{pressman2020, reenskaug1979}.

% ─────────────────────────────────────────────────────────────
\section{Restricciones y Fronteras de Diseño}
\label{sec:restricciones}

Para delimitar el alcance técnico y operativo del proyecto AUTECH, y garantizar su viabilidad
dentro de los tiempos de desarrollo del Trabajo Terminal, se establecen las siguientes
fronteras fundamentales:

\subsection{Restricciones Técnicas y de Hardware}
\label{ssec:restricciones_tecnicas}

\textbf{Segmentación de Plataforma:} El desarrollo se orienta de forma exclusiva al sistema
operativo Android, estableciendo la versión 9.0 (Pie, API~28) como requerimiento
mínimo~\cite{google2023android}. No se contempla compatibilidad con iOS debido a los costos
de licenciamiento asociados a la cuenta de desarrollador de Apple, además de la falta de
infraestructura para el empaquetado de recursos para iOS.

\textbf{Optimización Gráfica y Rendimiento:} El motor de renderizado Impeller, utilizado para
la manipulación fluida de autómatas en pantalla, aprovecha Vulkan como backend gráfico
principal cuando el dispositivo lo soporta, manteniéndose compatible mediante OpenGL~ES en
hardware que no cuente con dicha capacidad~\cite{flutter2024a}. El hardware mínimo requerido
consiste en un procesador de arquitectura de 64~bits y 4~GB de memoria RAM.

\textbf{Conectividad Híbrida:} El sistema operará bajo un modelo de tolerancia a
desconexiones. El motor de simulación y las lecciones teóricas funcionarán de manera autónoma
sin conexión a internet. No obstante, se requiere conectividad intermitente para la
autenticación de usuarios, la visualización de material audiovisual externo y la
sincronización de analíticas hacia el panel docente mediante Supabase~\cite{supabase2024}.

\textbf{Adaptabilidad de Interfaz:} El diseño garantizará la correcta visualización en
smartphones convencionales y tablets, priorizando una distribución de elementos que maximice
el área de trabajo del simulador en función del dispositivo del usuario
(responsive)~\cite{materialdesign2023}.

\subsection{Fronteras Didácticas y Operativas}
\label{ssec:fronteras_didacticas}

\textbf{Alcance Temático Institucional:} La base algorítmica y de contenidos se restringe
estrictamente a los tópicos del programa sintético de Teoría de la Computación
(ISC~2020)~\cite{ipn2020}. El avance lógico sigue la Jerarquía de Chomsky, Autómatas Finitos,
Autómatas de Pila y Máquinas de Turing, omitiendo deliberadamente temas de complejidad
computacional avanzada, como las clases P vs.\ NP, que exceden el propósito de esta
herramienta~\cite{hopcroft2007}, así como los lenguajes tipo~1 dentro de la Jerarquía de
Chomsky.

\textbf{Rol de Evaluación:} AUTECH está creada como una plataforma de refuerzo y evaluación
formativa. El software no pretende sustituir el proceso de evaluación sumativa oficial del
docente titular de la asignatura, sino tomar el rol de herramienta de apoyo a las clases
impartidas por el mismo~\cite{escom2016}.

\textbf{Estrategia Multimedia:} El contenido de video se define como recurso secundario. Las
lecciones están diseñadas para asimilarse íntegramente mediante textos, autómatas interactivos
e infografías, asegurando que la falta de ancho de banda no comprometa la integridad
didáctica de la unidad~\cite{crompton2018}.